\section{Análisis D6}

\subsection{Justificación del enfoque adaptativo}

Un enfoque predictivo presupone requisitos tempranamente estables, menor incertidumbre, cambios controlados y retroalimentación más tardía. En este caso, los requisitos pueden evolucionar, existe incertidumbre en los datos, riesgos y solución, y se requieren experimentación y revisión continua. Por ello, la adecuación del enfoque predictivo es media y la del enfoque adaptativo es alta.

El enfoque adaptativo no significa improvisación. Conserva una línea base de alcance y objetivos, pero permite ciclos incrementales con entregables, responsables, revisiones y puntos de decisión. Esta modalidad permite ajustar la planificación frente a la evidencia de calidad de datos, desempeño, riesgos, privacidad, sesgos y factibilidad de las alternativas.

\subsection{Criterios completos de comparación}

\begin{description}[style=nextline]
  \item[Requisitos] En un enfoque predictivo son estables tempranamente; en uno adaptativo pueden evolucionar. En este caso, pueden evolucionar.
  \item[Incertidumbre] Es baja en el enfoque predictivo y media o alta en el adaptativo. El caso presenta incertidumbre alta en datos, riesgos y solución.
  \item[Cambios] El enfoque predictivo los minimiza o controla; el adaptativo los incorpora de forma planificada. En este proyecto son probables.
  \item[Retroalimentación] Es más tardía en el enfoque predictivo y frecuente en el adaptativo. El caso necesita retroalimentación frecuente.
  \item[Datos y modelos] El enfoque predictivo requiere mayor previsibilidad; el adaptativo permite experimentar. El caso presenta incertidumbre en esta materia.
  \item[Riesgos] En el enfoque predictivo se planifican principalmente al inicio; en el adaptativo se revisan durante los ciclos. El caso requiere revisión continua.
  \item[Adecuación] La adecuación predictiva es media y la adaptativa es alta; por ello se selecciona el enfoque adaptativo.
\end{description}

El efecto transversal es claro: la EDT y la carta Gantt deben considerar revisiones y capacidad de ajuste; finanzas debe presupuestar iteración y contingencia; riesgos debe actualizarse al aparecer nueva evidencia; y la evaluación ética debe acompañar el ciclo completo.
